% ============================================================
% 2. Literature Review
% ============================================================
\section{Related Work}
\label{sec:lit}

This paper sits at the intersection of three literatures: autonomous agents
in scientific and engineering design, classical optimization baselines and
their known boundaries, and how delegated decisions are checked. We
review each in turn and then state the gap.

\subsection{Autonomous Agents in Scientific and Engineering Design}

The past three years have seen LLM agents deployed across the design chain,
one slice at a time. Coscientist automated end-to-end chemistry workflows
 (planning, coding, instrument control) in a single agent
\cite{boiko2023coscientist}; ChemCrow augmented an LLM with chemistry tools
for molecule design and synthesis planning \cite{bran2024chemcrow}; the A-Lab
closed the loop from synthesis recipe to robotic execution to X-ray
characterization for materials discovery \cite{szymanski2023alab}. In the
battery domain specifically, the lineage is young but active: LLM-assisted
battery research has been surveyed comprehensively
\cite{batteryLLMreview2025}; foundation models target battery analysis
\cite{lipm2025}; and multi-agent chemist systems extend the approach to
material prediction and exploration \cite{chatmat2026, matsciagent2025,
matclaw2026}, and further agent families span hierarchical multi-agent
materials discovery \cite{rothfarb2025hierarchical}, drug piloting and
RAG-enhanced collaborative drug discovery \cite{drugpilot2025, cladd2025},
and AMS circuit-design optimization \cite{heart2025}. Closest to our
setting, Battery-Sim-Agent performs
\emph{inverse} battery parameter estimation with an LLM agent in a
simulator-in-the-loop configuration, reporting 67--95\% curve-matching error
reduction against classical black-box optimizers across 200 synthetic
calibration tasks \cite{chen2026battery}. Adjacent engineering domains
replicate the pattern: AnalogAgent automates analog circuit design
\cite{analogagent2026}; GENIUS orchestrates autonomous design and execution
of simulation runs \cite{genius2025}.

Two structural observations hold across this entire body of work. First,
every system covers a slice of a design chain, not a chain: material
discovery stops at the material, circuit automation stops at the circuit,
inverse estimation stops at parameters. None spans a full set of coupled
design variables (from chemistry through formulation through architecture) and
none therefore faces the cross-scale diagnosis problem that a full chain
poses. Second, the judgment in each system is asserted by the agent itself.
Whether a synthesis ``worked'' or a parameter fit ``converged'' is reported in
natural language; there is no pre-registered target against which a
code-level verifier judges the claim. Benchmark evaluations of these systems
reflect the same design: DiscoveryBench, ScienceAgentBench, MLAgentBench, and
RE-Bench measure task success and capability \cite{discoverybench2024,
scienceagentbench2024, mlagentbench2023, rebench2024}; FIRE-Bench adds an
honesty layer for scientific rediscovery claims \cite{firebench2026}. None
treats \emph{whether the agent reports against the metric it was given} as a
first-order experimental
variable.

\subsection{Classical Optimization Baselines and Their Known Boundaries}

Bayesian optimization (BO) is the standard strong baseline against which LLM
agent systems are measured, and the comparison record already contains an
important calibration. Battery-Sim-Agent benchmarks its agent against
Meta's Ax BO platform (Sobol initialization, GPEI with LogNEI acquisition,
fixed seed) and reports not only the headline reduction but a
calibrated boundary: in ``low-headroom'' regimes where literature default parameters
already sit close to the target, BO remains competitive and can even
outperform the agent, and CMA-ES fails to converge on their parameter
estimation tasks altogether \cite{chen2026battery}. We adopt the identical BO
configuration (same platform, same strategy family, same seed) so that our
baseline inherits this published calibration.

The deeper boundary, which no agent study has isolated experimentally, is
structural rather than statistical: BO searches a fixed parameter space. A
classical optimizer can move thickness and porosity, but it cannot invent an
electrolyte formulation, propose a coating, or switch a cathode system. When
the binding constraint on a design target is a material property---a
voltage plateau beyond the chemistry's polarization limit, an SEI growth
rate that only a coating suppresses: the optimum is not in the reachable
set, and no acquisition function can find it. The battery modeling
literature supplies the physics that makes this boundary concrete: plating
onset is governed by the local overpotential balance
\cite{arora1999plating, ren2018plating}; SEI growth by reaction kinetics
\cite{peled1979sei}; thermal runaway by the coupled decomposition chemistry
\cite{wang2012thermalrunaway}. Each of these mechanisms admits variables at
scales a structural optimizer cannot touch. Our experimental design makes
this boundary an observable: three of our eight targets are solvable
within the structural space, and five are not.

\subsection{Agent Reliability and the Case for External Verification}

The central obstacle to trusting an agentic design system is not model
capability but measurement: an agent that both executes and evaluates can
shift the criterion it is measured against. The AI community has begun to
quantify this. Honesty benchmarks establish that LLMs systematically
overstate what they know \cite{behonest2024, honesty2023alignment};
FIRE-Bench shows that even frontier scientific agents fail most reported
reconstructions, and that their self-reported completion is not a reliable
signal of actual convergence \cite{firebench2026}. Evaluation practice has
responded with two remedies. The first is \emph{soft judging}: LLM-as-judge
and reward-model reviewers applied to agent outputs. The second is
\emph{hard judging}: deterministic, code-based evaluation of claims against
ground truth (as in unit-test-style benchmark harnesses). Hard judging is
the technically reliable option, but in existing systems it evaluates
\textbf{the benchmark}, not the agent's own workflow; the design loop
itself is never forced under a deterministic evaluator. The management literature arrives at the same
conclusion from the opposite direction: delegated agency theory
\cite{delegatedagency2026} identifies how well the target is specified and
whether the workflow's decisions are visible as the design dimensions that discipline delegated
action; empirical work on AI-assisted decisions shows algorithmic
recommendations reshape human override behavior
\cite{nikpayam2024aiassisted}; Bayesian delegation policies allocate
decision authority sequentially \cite{dixon2026bayesiandelegation}; GAIA
articulates information-gated progression for high-stakes LLM agents
\cite{zhao2025gaia}; and qualitative work warns that autonomy must be
designed into systems rather than assumed
\cite{dancingalgorithm2025}. What none of this literature contains is an
\emph{engineering-level} instantiation: a concrete design environment in
which the evaluator is code, the evaluation target is pre-registered
before work begins, and the fidelity of that evaluation is itself an
experimental variable.

Finally, we note the disciplinary bridge this paper crosses. The battery
prognostics community rooted in grey system theory (originating in the
management science tradition \cite{deng1989grey}) has long modeled battery
degradation and state-of-charge as forecasting problems
\cite{li2025greyiterative, li2024hybridgrey, sapnken2025soh,
sapnken2024pemfc, sapnken2026greyneural, xu2025greykalman}. That lineage
treats the battery as an object to be \emph{predicted}. This paper treats it
as an object to be \emph{designed by an agent that follows a workflow},
extending the same lens from prognostics to how design decisions are checked.

\subsection{The Gap and Research Questions}

Three gaps follow. (1) No published system spans the full chain of cell
design variables (electrolyte, separator, electrode, architecture,
thermal) in one autonomous loop. (2) No experimental comparison isolates
the two failure modes of ungoverned delegation: the solution-space failure of
classical optimizers and the measurement drift of workflow-free agents. (3)
No study treats the evaluator (pre-registered targets,
code-code-computed verdicts, evidence-chain audits) as an experimental
variable whose removal can be measured. The paper is designed to isolate
both failure modes: the solution-space failure and the measurement drift.

We formalize three research
questions:

\begin{itemize}
\item[\textbf{RQ1}]
\emph{The design workflow performance.} Does a fixed design workflow raise target
attainment relative to a workflow-free agent and a Bayesian optimizer under
identical resources, and at what resource (token and message) cost?
\item[\textbf{RQ2}]
\emph{Mechanism attribution.} Which workflow rules carry the effect,
and in which task regimes? We ablate each component separately and report
its marginal contribution, including null results.
\item[\textbf{RQ3}]
\emph{Can the reported measurement be trusted?} Does code-level evaluation close the
judgment-integrity gap (measured as evidence completeness, criterion
fidelity, and the absence of self-set thresholds) that the workflow-free
agent exhibits?
\end{itemize}
